\documentclass[conference]{IEEEtran}
\IEEEoverridecommandlockouts
% The preceding line is only needed to identify funding in the first footnote. If that is unneeded, please comment it out.
%Template version as of 6/27/2024

\usepackage{cite}
\usepackage{amsmath,amssymb,amsfonts}
\usepackage{algorithmic}
\usepackage{graphicx}
\usepackage{textcomp}
\usepackage{xcolor}
\usepackage{url}
\usepackage{hyperref}
\usepackage{ctable}
\usepackage{subcaption}
\usepackage{tabto}    
\usepackage{array}
\usepackage{multirow}
\usepackage{float}
\usepackage{dblfloatfix}
\usepackage{standalone}

\usepackage{pdfpages}



\def\BibTeX{{\rm B\kern-.05em{\sc i\kern-.025em b}\kern-.08em
    T\kern-.1667em\lower.7ex\hbox{E}\kern-.125emX}}

\newcommand{\linebreakand}{%
  \end{@IEEEauthorhalign}
  \hfill\mbox{}\par
  \mbox{}\hfill\begin{@IEEEauthorhalign}
}

\begin{document}

\title{Uncertainty-Aware Kolmogorov-Arnold Networks for Improved Medical Out-of-distribution Detection\\
}

\author{\IEEEauthorblockN{Konstantinos Konstantinidis}
\IEEEauthorblockA{\textit{School of Informatics} \\
\textit{Aristotle University of Thessaloniki}\\
Thessaloniki, Greece \\
\href{mailto:kdkonstanti@csd.auth.gr}{kdkonstanti@csd.auth.gr}}
\and
\IEEEauthorblockN{Dimitrios Spanos}
\IEEEauthorblockA{\textit{School of Informatics} \\
\textit{Aristotle University of Thessaloniki}\\
Thessaloniki, Greece \\
\href{mailto:dvspanos@csd.auth.gr}{dvspanos@csd.auth.gr}}
\and
\IEEEauthorblockN{Nikolaos Passalis}
\IEEEauthorblockA{\textit{School of Chemical Engineering} \\
\textit{Aristotle University of Thessaloniki}\\
Thessaloniki, Greece \\
\href{mailto:passalis@auth.gr}{passalis@auth.gr}}
\and[\hfill\mbox{}\par\mbox{}\hfill]
\IEEEauthorblockN{Anastasios Tefas}
\IEEEauthorblockA{\textit{School of Informatics} \\
\textit{Aristotle University of Thessaloniki}\\
Thessaloniki, Greece \\
\href{mailto:tefas@csd.auth.gr}{tefas@csd.auth.gr}}
}

\maketitle

\begin{abstract}
While state-of-the-art deep learning models have achieved remarkable performance across various domains, their widespread adoption in safety-critical applications remains constrained by significant limitations in uncertainty estimation and inherent overconfidence. These shortcomings pose substantial risks in domains requiring reliable confidence measures. Kolmogorov-Arnold Networks (KANs) are a compelling alternative, offering a unique combination of interpretability, lightweight architecture with their learnable activation functions. In this work, we propose a KAN architecture that unifies the structural advantages of KANs with the localized uncertainty modeling of Radial Basis Functions (RBFs), forming a single-forward-pass model optimized through a proposed composite loss for joint prediction and confidence estimation. Through comprehensive benchmarking on medical datasets, we demonstrate that our proposed Uncertainty-Aware Kolmogorov-Arnold Network (UA-KAN) achieves state-of-the-art OoD detection performance, significantly outperforming established baselines, while maintaining competitive accuracy and favorable training times. Code is available at: \href{https://github.com/konthis/Uncertainty-Aware-KAN-for-OoD-detection}{GitHub}. 
\end{abstract}


\begin{IEEEkeywords}
trustworthy artificial intelligence, deep learning, kolmogorov-arnold networks, uncertainty estimation, out-of-distribution detection.
\end{IEEEkeywords}


\section{Introduction}
Deep Learning systems are rapidly evolving across diverse domains, from critical applications such as financial fraud detection and autonomous driving to industrial automation and everyday personal solutions \cite{10879299} \cite{YAO2024130419} \cite{Ahmed2025}. Although the goal is always to develop high-performance and accurate systems, not all problems are the same. In safety-critical domains, such as medical diagnosis, finance, etc., high trustworthiness is an absolute necessity \cite{ZOU2023100002}, as all low trustworthiness systems will never reach common use, because the consequences can be high-cost. One big contributor to trustworthiness is reliable uncertainty estimation \cite{YAO2024130419}. A system with reliable uncertainty estimation, quantifies confidence, detects Out-of-Distribution (OoD) data reliably, and can trigger human intervention when it encounters ambiguous or unfamiliar scenarios. Developing accurate systems must come to par with high uncertainty estimation power, and, ultimately, integrating trustworthy uncertainty estimations into deep learning systems paves the way for their responsible deployment in critical domains, balancing performance with reliability.

Deep learning enables the construction of computational models with multiple processing layers that learn hierarchical data representations at increasing levels of abstraction \cite{LeCun2015}. Although Deep Learning systems solve a wide range of tasks, traditional architectures, such as Multi-Layered Perceptrons (MLPs), suffer from poor uncertainty estimation and  overconfidence \cite{mucsányi2023trustworthymachinelearning}, making them untrustworthy, hindering their widespread adoption in safety-critical applications, such as medical diagnosis, autonomous driving, financial decisions, etc. This is a result of the insensitivity of these systems in anomalies found in input data such as in one-pixel attacks \cite{su2019one}, creating overconfident predictions, even on OoD data. While methods like Monte Carlo Dropout \cite{10401930} and Deep Ensembles \cite{fort2020deepensembleslosslandscape} have been proposed to improve uncertainty estimation, they often come with significant computational overhead. Recently, Radial Basis Functions (RBF) architectures have emerged as a promising alternative to provide better uncertainty estimates \cite{11087916} in a lightweight manner. 

Kolmogorov–Arnold Networks (KANs) offer a structural advantage for uncertainty estimation because they represent functions through learned univariate transformations composed along low dimensional interactions \cite{liu2025kankolmogorovarnoldnetworks}. This architecture encourages local adaptivity since each edge learns its own spline function rather than relying on a fixed global nonlinearity. As a result, the model can represent sharp variations, flat regions, and boundary behaviors of the data manifold in a more controlled manner. Conventional MLPs often extrapolate confidently in such regions because their representations are dominated by global linear transformations combined with point-wise activations. The functional flexibility of KANs can therefore provide a way for expressing heteroscedasticity and local feature dispersion. Moreover, RBF architectures model uncertainty through distance metrics in their representation space, where the confidence of a prediction is inherently linked to its proximity from a learned centroid \cite{SPANOS2025130954}. However, there are challenges in getting reliable uncertainty estimates from KANs, as their additive functional decomposition can amplify the contributions of multiple univariate components, leading to large logit magnitudes and overconfident predictions even in regions where the input lies far from the training distribution.


These observations lead to a key research question: Can an uncertainty-aware Kolmogorov–Arnold Network improve the performance of traditional deep learning models in uncertainty estimation and out-of-distribution (OoD) detection while maintaining competitive in-domain accuracy? In this paper, we address this question by proposing an uncertainty-aware Kolmogorov–Arnold Network (UA-KAN) that replaces the standard B-spline parameterization with fully learnable radial basis function (RBF) activations. Unlike spline-based edges, RBF units exhibit localized and isotropic responses that decay smoothly with distance from the training distribution, yielding distance-aware representations that naturally reflect epistemic uncertainty and reduce overconfident extrapolation. Building on this architectural choice, we introduce an optimization framework that jointly promotes classification accuracy and uncertainty calibration through a composite loss formulation, and we further incorporate logit normalization into the KAN framework to explicitly regulate logit scale and align output confidence with the localized structure of RBF-based representations. Our main contributions are as follows: (1) an uncertainty-aware Kolmogorov–Arnold Network model featuring fully learnable RBF activation functions and logit-normalized predictions, and (2) a custom loss function that enforces both discriminative performance and calibrated uncertainty. Finally, we provide a comprehensive evaluation comparing UA-KAN against established baselines, including MLPs, Deep Ensembles \cite{lakshminarayanan2017simplescalablepredictiveuncertainty}, DUQ \cite{vanamersfoort2020uncertaintyestimationusingsingle}, and existing KAN variants, demonstrating that our method consistently achieves superior OoD detection performance in terms of AUROC while retaining competitive in-domain accuracy.

The rest of the paper is structured as follows. In Section 2 we present an overview of related work in the field. Following that, in Section 3, the proposed method and an analysis of the architecture are presented. The experimental setup and the experimental evaluation, are provided in Section 4. Finally, in Section 5 conclusions are drawn.

\section{Related Work}


Uncertainty estimation in deep learning has been approached through several complementary paradigms. Bayesian Neural Networks (BNNs) \cite{Goan_2020} offer a principled probabilistic framework via Bayes' Theorem, though approximate inference limits their scalability \cite{gal2016dropoutbayesianapproximationrepresenting}. MC Dropout \cite{gal2016dropoutbayesianapproximationrepresenting} approximates Bayesian inference at test time through stochastic forward passes, while Deep Ensembles \cite{lakshminarayanan2017simplescalablepredictiveuncertainty} boost robustness by aggregating predictions across independently trained models at a proportionally higher cost. DUQ \cite{vanamersfoort2020uncertaintyestimationusingsingle} uses an RBF network with a gradient penalty for single-pass uncertainty, though it often trades accuracy for calibration. Energy-based OoD detection \cite{liu2020energy} derives a free energy score from output logits, offering more principled in/OoD separation than softmax confidence \cite{10.1007/978-3-031-78201-5_4}. NeuralFP \cite{9412489} performs OoD detection via layer-wise autoencoder fingerprinting, while CTR \cite{9956206} addresses the broader Open-World Recognition setting through contrastive learning.

Tabular data dominates clinical applications, with XGBoost \cite{chen2016xgboost}, TabNet \cite{arik2021tabnet}, and FT-Transformer \cite{gorishniy2021revisiting} representing strong discriminative baselines; however, none is designed for uncertainty quantification or OoD detection. TabPFN \cite{hollmann2023tabpfn} is a prior-fitted Transformer that solves small tabular classification problems via in-context learning, processing the entire training set at inference time to achieve strong performance without task-specific training; however, this yields inference costs that scale quadratically with training set size. On the KAN side, the original architecture \cite{liu2025kankolmogorovarnoldnetworks} uses B-splines, which are inefficient on GPU; Convolutional KANs \cite{bodner2025convolutionalkolmogorovarnoldnetworks} extend this to vision tasks, while FastKAN \cite{li2024kolmogorovarnoldnetworksradialbasis} replaces B-splines with RBFs for improved efficiency. Unlike these variants, our method integrates per-edge uncertainty directly into the KAN's functional structure and couples it with logit normalization, jointly addressing in-domain accuracy and OoD detection in a single model.


\begin{figure*}[h]
  \centering
  \includegraphics[width=1\textwidth, angle=0]{figures/ua_kan_architecture.pdf}
  \caption{UA-KAN architecture with per-layer uncertainty estimation.}
  \label{fig:ua_kan_arch}
\end{figure*}

\section{Proposed Method}

This section details the contributions of our work. We establish the necessary background on KANs and the task of OoD detection. Moreover, we provide a theoretical motivation for using RBF-based activations over B-splines and an analysis of the proposed architecture. Finally, we introduce and motivate the training objective, including the rationale for the formulation of the proposed loss.

\subsection{Background \& Notation}

\noindent \textbf{Kolmogorov-Arnold Networks} Let $\mathbf{x} = [x_1, x_2, \dots, x_n] \in \mathbb{R}$, where $n$ is the input dimension, be the input vector of a KAN and $y\in \mathbb{R}^C$ be the target, where $C$ denotes the number of classes. The Kolmogorov–Arnold representation theorem states that any continuous multivariate function $f(x_1, x_2, \ldots, x_n)$ can be represented as a finite superposition of continuous univariate functions and addition. Specifically, there exist continuous functions $\Phi = \{\phi_{q,p}\},\,  p = 1,2,\ldots,n_{in}, \, q = 1,2,\ldots,n_{out}$ such that:

\begin{equation} \label{eq:kantheor1}
f(x_1, x_2, \ldots, x_n) = \sum_{q=1}^{2n+1}\Phi_q(\sum_{p=1}^{n}\phi_{q,p}(x_p)),
\end{equation}
where $n_{in}$ is the input data dimension, $n_{out}$ is the output data dimension, $\phi_{q,p}$ are univariate functions that map each input variable $x_p$, such that $\phi_{q,p}:[0,1]\xrightarrow{}\mathbb{R}$, and $\Phi_q:\mathbb{R}\xrightarrow{}\mathbb{R}$ are also univariate functions.
%https://arxiv.org/abs/2404.19756
Equation \ref{eq:kantheor1} represents a two-layer KAN, where the hidden layer consists of a relatively small number of terms $(2n+1)$. However, it has been shown \cite{liu2025kankolmogorovarnoldnetworks} that this representation can be generalized to architectures of arbitrary depth and width, thereby enabling the construction of deeper and more expressive networks based on the same theoretical foundation.

To give a general KAN composition in mathematical form, given an input $\mathbf{x}\in \mathbb{R}$ and L layers, the output is defined as:
\begin{equation} \label{eq:kancomposition}
KAN(\mathbf{x}) = (\mathbf{\Phi}_{L-1}\circ\mathbf{\Phi}_{L-2}\circ\dots\circ\mathbf{\Phi}_{0})\,\mathbf{x}
\end{equation}
From the perspective of a neural network, the goal is to learn a target function $\mathbf{\hat{y}} = f(\mathbf{x})$ from data by adjusting the parameters of the univariate functions $\Phi_q$ and $\phi_{q,p}$ at each layer, where $\hat{\mathbf{y}}\in \mathbb{R}^C$ is the pre-softmax output of the model.


The KANs reframes this multivariate learning problem into discovering a composition of univariate functions, as defined by equation \ref{eq:kantheor1}. In practice, these unknown univariate functions $\Phi_q$ and $\phi_{q,p}$, are represented by parameterized models. Specifically, each function is implemented as a linear combination of basis functions with trainable coefficients. The original KAN architecture utilizes B-spline bases, whereas our proposed UA-KAN adopt RBF kernels for this purpose.

% https://arxiv.org/pdf/2404.19756


In this setting, each univariate function $\phi(x)$ is implemented as a linear combination of two components: a fixed basis function $b(x)$ and a learnable B-spline function:
\begin{equation} \label{eq:kanimplement}
\phi(x) = w_b \, b(x) + w_s \, spline(x),
\end{equation}
where $w_b$ and $w_s$ are trainable scaling coefficients and $b(x) = SiLU(x)$.

% https://arxiv.org/pdf/2404.19756

\noindent \textbf{Out-of-distribution Detection} For reliable real-world deployment, a classifier must not only be accurate on known in-domain data, but also correctly identify and reject OoD inputs. Formally, given a classifier $f$, trained on a set of in-domain data, the goal of an OoD detector is to identify whether a test sample originates from a different distribution. 

A standard detection mechanism uses a scoring function ($S$) and a threshold $\lambda$, classifying samples with scores above $\lambda$ as in-domain (ID) and those below as OoD:
\begin{equation}
\label{eq:oodthreshold}
G(x) = \begin{cases}
ID &\text{if $S(x) \ge \lambda$}\\
OoD &\text{if $S(x) < \lambda$}
\end{cases}
\end{equation}
%%%%%%%%%%%%%%%%
Performance is measured using the Area Under the Receiver Operating Characteristic curve (AUROC) \cite{10.1007/978-3-031-78201-5_4}, a threshold-independent metric that quantifies the ability of the scoring function $S$ to separate in-distribution from OoD samples across all possible thresholds $\lambda$.
%%%%%%%%%%%%%%%

\subsection{Proposed UA-KAN architecture}

A key motivation for our architectural choice stems from a structural analogy with DUQ \cite{vanamersfoort2020uncertaintyestimationusingsingle}, which applies RBF-based distance scoring exclusively at the final output layer. In contrast, when every learnable weight in a KAN is replaced by a 1D RBF function, each edge becomes a miniature distance-aware unit, naturally tracking how familiar the input is at every step of the computation. For a given edge with learned centers $\mu_i$, bandwidth $\gamma$, and scalar input $x$, the output prediction is computed as a weighted RBF expansion:

\begin{equation} \label{eq:edgepred}
y = \sum_i c_i \exp\!\left(-\gamma(x - \mu_i)^2\right),
\end{equation}

where $c_i$ are learned weights. Simultaneously, the epistemic uncertainty at that edge is quantified as:

\begin{equation} \label{eq:edgeuncert}
U_{\text{edge}} = 1 - \max_i \left(\exp\!\left(-\gamma(x - \mu_i)^2\right)\right),
\end{equation}

which approaches unity when $x$ lies far from all learned centers, indicating that the input falls outside the support of the training distribution at that edge. These per-edge uncertainty signals are aggregated across all layers to form a global confidence estimate passed to the OoD detector, as illustrated in Fig.~\ref{fig:ua_kan_arch}. This mechanism fundamentally distinguishes UA-KAN from DUQ and FastKAN \cite{li2024kolmogorovarnoldnetworksradialbasis}: rather than appending uncertainty as an output-layer post-processing step, it is embedded in the functional structure of every network edge.


RBF-based KANs provide more reliable uncertainty estimates than B-spline KANs because their basis functions have an inherently localized and isotropic response. Each RBF unit captures information in a strictly bounded neighborhood of the input space, causing activations to decay smoothly as inputs move away from regions supported by the training data. This natural “distance-aware” behavior yields sharper epistemic uncertainty: when the model encounters OoD samples, the RBF activations gradually converge to zero, limiting extrapolation and preventing spurious high-confidence predictions. In contrast, B-splines rely on piecewise polynomials with compact but axis-aligned support; outside their knot intervals, they often extrapolate linearly or remain constant, which weakens their ability to reflect how far an input is from known data manifolds. 

Moreover, RBF-based KANs create a smoother, more stable geometry in representation space, which directly benefits downstream OoD metrics. Because RBF expansions are radially symmetric, they approximate likelihood-like behavior and distribute representational density according to training data concentration. This structure aligns naturally with uncertainty measures such as entropy, energy scores, or Dirichlet concentration, enabling the model to separate in-domain and OoD inputs more distinctly. B-spline KANs, while expressive, typically exhibit anisotropic and grid-dependent behavior: their uncertainty depends heavily on knot placement and coordinate axes, making them more susceptible to overconfident extrapolation in sparse regions. As a result, RBF-based variants tend to yield clearer decision boundaries, more calibrated confidence, and higher robustness in OoD detection tasks.

Our proposed UA-KAN replaces the B-spline representation of the univariate functions $\phi_{q,p}$ with a formulation based on Radial Basis Functions (RBFs), while addressing the training-time inefficiency of the original KAN's B-spline parameterization. For a given input 
$\mathbf{x}$, the function is approximated as:
% Specifically, referring to FastKAN paper \cite{li2024kolmogorovarnoldnetworksradialbasis} we  
\begin{equation} \label{eq:rbfkan}
RBF(\mathbf{x}) = \sum^N_{i=1} w_i\phi_{RBF}(||\mathbf{x}-\mathbf{c}_i||_2),
\end{equation}
where $|| \cdot ||_2$ is the $L_2$ norm, $w_i$ are learnable weights, $c_i$ the RBF centers, and $\phi_{RBF}$ is the Gaussian function:   
\begin{equation} \label{eq:rbfeq}
\phi_{RBF}(r) = exp(\frac{-r^2}{2h^2}),
\end{equation}
where $r$ is the radial distance and $h$ is the width of the RBF. Furthermore, to improve uncertainty estimation, we propose the RBF-SiLU (RS) as the residual function for the basis component. This function replaces the standard SiLU in Eq. (\ref{eq:kanimplement}) as follows:

% https://arxiv.org/pdf/2405.06721



\begin{equation} \label{eq:rbfsilu}
RS(x) = xexp(-\gamma x^2),
\end{equation}
where $\gamma$ contributes to the RBF width. Our implementation maintains the original KAN's design principle: each univariate function $\phi(x)$ is computed as a weighted combination of a baseline function (RS), and the RBF function:

\begin{equation} \label{eq:ourkanimplement}
\phi(x) = w_bRS(x) + w_sRBF(x).
\end{equation}
where $w_b$ and $w_s$ are trainable scaling coefficients. To increase the flexibility of the learning process, we parameterize the width $h$ of the RBF kernel, Eq. \ref{eq:rbfeq}, as a trainable parameter, allowing the model to learn the optimal locality and scale for each activation function.

%%%%%%%%%%%%%%%%%%%%%%%%%%%%%%%%%%%%%%

\subsection{Per-Layer Uncertainty Estimation}

The RBF basis activations at each layer provide a natural uncertainty signal. For a given layer with input $\mathbf{x}$, let $\boldsymbol{\phi} \in \mathbb{R}^{n_{in} \times N}$ denote the RBF basis activations, where $n_{in}$ is the number of input features and $N$ is the number of RBF centers. Each row $\boldsymbol{\phi}_j \in \mathbb{R}^N$ is normalized across centers to form a distribution:

\begin{equation} \label{eq:rbfnorm}
\tilde{\boldsymbol{\phi}}_j = \frac{\boldsymbol{\phi}_j}{\sum_{k=1}^{N} \phi_{j,k} + \epsilon},
\end{equation}

and the layer-level uncertainty is the mean entropy of these per-feature distributions:

\begin{equation} \label{eq:layerentropy}
H_l = \frac{1}{n_{in}} \sum_{j=1}^{n_{in}} \left( -\sum_{k=1}^{N} \tilde{\phi}_{j,k} \log\left(\tilde{\phi}_{j,k} + \epsilon\right) \right).
\end{equation}

When the input is in-distribution, activations concentrate near learned centers, yielding low entropy. When the input is OoD, no center is strongly activated and the distribution becomes uniform, yielding high entropy.

We evaluate two variants. In the simple variant, the OoD score is derived solely from $H_L$ at the final layer. In the weighted variant, layer uncertainties are aggregated as:

\begin{equation} \label{eq:weighteduq}
U = \sum_{l=1}^{L} w_l \, H_l,
\end{equation}

where weights $w_l$ are either uniform ($w_l = 1/L$) or set post-training to the variance of layer activations over $\mathcal{D}_{train}$:

\begin{equation} \label{eq:layerweights}
w_l = \mathrm{Var}_{x \in \mathcal{D}_{train}}\!\left[\mathbf{a}^{(l)}(x)\right],
\end{equation}

assigning higher influence to layers whose activations vary most across the training distribution. The score $U$ is then passed to the OoD detector, as shown in Fig.~\ref{fig:ua_kan_arch}.


%%%%%%%%%%%%%%%%%%%%%%%%%%%%%%%%%%%%%%%%

\subsection{Proposed UA-KAN Objective}

Following the introduction of the proposed UA-KAN architecture, we motivate the use of logit normalization within the KAN framework. Neural networks trained with cross-entropy tend to produce logits of unbounded magnitude, which inflate softmax confidence and cause overconfident predictions on OoD inputs \cite{wei2022mitigating}. Logit normalization addresses this by projecting the logit vector onto the unit sphere before computing the softmax, explicitly decoupling the direction of $\mathbf{f}$ (which encodes class preference) from its magnitude (which inflates confidence), and has been shown to significantly improve OoD detection in MLP-based models \cite{wei2022mitigating}. Although logit normalization has been shown to improve calibration in MLP-based models \cite{wei2022mitigating}, its effect is limited by the globally entangled representations produced by such architectures, where confidence regulation at the output layer is weakly coupled to the internal feature structure. In contrast, KANs rely on an additive functional decomposition that yields more structured and disentangled representations. When trained with cross-entropy alone, this structure can still give rise to excessively large logits, leading to overconfident predictions that do not faithfully reflect the underlying evidence. Logit normalization mitigates this issue by explicitly constraining logit scale while retaining relative class preferences, making it intrinsically more compatible with the KAN decision process than with standard MLPs.

This advantage becomes more pronounced in RBF-based KANs with radial receptive fields, where each component encodes localized responses and distance to the training distribution. In this setting, output confidence should emerge from the relative contributions of these localized functions rather than from absolute activation magnitude. Logit normalization enforces this behavior at the prediction stage, effectively coupling confidence regulation with the locality induced by the proposed architecture. As a result, the combination of KANs and logit normalization yields a tighter alignment between representation, decision making, and uncertainty than is typically achievable in MLPs, providing a clear motivation for the loss formulation introduced next.

Our model is optimized using a custom loss function designed to jointly address discriminative accuracy and uncertainty evaluation. This function is defined as a weighted linear combination of Cross Entropy (CE) and Logit Norm (LN) loss, which mitigates overconfidence and we empirically found to yield superior final performance compared to using either loss in isolation:

\begin{equation} \label{eq:Categoricalcrossentr}
\mathcal{L}_{CE} = - \sum_{c=1}^{C}y_{c}\,log( p_{c}),
\end{equation}
where C is the number of classes and $\hat y$ the predicted class and:

\begin{equation} \label{eq:logitnormloss}
\mathcal{L}_{LN} = - log\frac{e^{(\frac{f_i}{\tau||\mathbf{f}||_2})}}{\sum_{j=1}^{C}e^{(\frac{f_j}{\tau||\mathbf{f}||_2})}},
\end{equation}
where $f_i$ is the i-th logit of the model, $\tau$ is the temperature parameter of the logits and $||\mathbf{x}||_2$ is the $L_2$ norm of the inputs. So, the proposed loss is formed as:
\begin{equation} \label{eq:proposedloss}
\mathcal{L}_{proposed} = \lambda \, \mathcal{L}_{CE} + (1-\lambda) \, \mathcal{L}_{LN},
\end{equation}
where $\lambda$ being the mixing coefficient.




\section{Experimental Evaluation}

This section presents the evaluation of our proposed UA-KAN method against established uncertainty quantification and OoD detection baselines. We introduce the datasets used for in-domain training and testing, along with the OoD datasets. Then we detail the experimental setup, including the specific configurations and hyperparameters for all models. The results are sectioned on multi-class and binary-class classification performance, computational efficiency, and an analysis on our proposed loss function. 


\input{tables/ambrosia3class.tex}

\input{tables/ambrosia2class.tex}





\subsection{Datasets}

\noindent \textbf{SEPSIS} We evaluate our models on the clinical SEPSIS dataset (Version 1) \cite{11087847}, specifically designed for sepsis detection. Each sample is annotated by clinical experts based on three key inflammatory biomarkers available early in a patient's admission: C-reactive protein (CRP), Procalcitonin (PCT), and Interleukin-6 (IL-6). The original multi-class task categorizes each sample into one of three clinical states: Non-Infectious Systemic Inflammatory Response Syndrome (SIRS), Sepsis, and Septic Shock. We derive a binary task by merging the Sepsis and Septic Shock categories into a single ``Infection'' positive class, distinguished from the Non-Infectious SIRS negative class, creating a clinically relevant task of identifying any infectious systemic inflammatory response. To evaluate uncertainty estimation and OoD detection, we employ three standard benchmark datasets from the UCI Machine Learning Repository via the scikit-learn library \cite{pedregosa2011scikit}: Wine \cite{wine_109}, Iris \cite{iris_53}, and Breast Cancer (Diagnostic) \cite{breast_cancer_wisconsin_(diagnostic)_17}. Furthermore, we create a ``Noised SEPSIS'' dataset by applying heavy Gaussian noise to simulate an OoD dataset.

\noindent \textbf{Heart Disease} We further evaluate on the Heart Disease dataset \cite{heart_disease_kaggle}, a multi-source compilation of clinical records described by 13 attributes, including age, chest pain type, resting blood pressure, cholesterol, and ST-segment parameters. The task comprises 5 clinically interpretable severity levels ranging from no disease to critical heart disease. OoD detection uses the Breast Cancer dataset \cite{breast_cancer_wisconsin_(diagnostic)_17} and a Gaussian-noised version of the training data.

All features were standardized to zero mean and unit variance prior to training.


\subsection{Experimental Setup}
All models were implemented in PyTorch with a learning-rate plateau scheduler (patience 5) and evaluated over 5 independent random seeds to quantify variance.

\noindent \textbf{Baselines} MLP and Deep Ensemble \cite{lakshminarayanan2017simplescalablepredictiveuncertainty} use architectures $[3,32,16,c]$ / $[13,128,64,5]$, trained with SGD (learning rate $0.1$, batch size $16$); the ensemble aggregates 5 independent MLP outputs. Energy MLP \cite{liu2020energy} shares this configuration, replacing softmax confidence with $E(\mathbf{x}) = -\log\sum_c e^{f_c(\mathbf{x})}$. MC Dropout \cite{gal2016dropoutbayesianapproximationrepresenting} uses $[3,64,32,c]$ / $[13,64,32,5]$ with dropout $p=0.2$, SGD (learning rate $0.1$), and 20 stochastic forward passes at inference. DUQ \cite{vanamersfoort2020uncertaintyestimationusingsingle} uses $[3,32,16,c]$ / $[13,64,32,5]$, SGD (learning rate $0.1$, 120 epochs), gradient penalty $\lambda=0.25$, bandwidth $\sigma=1$. FT-Transformer \cite{gorishniy2021revisiting} is trained with Adam (learning rate $10^{-4}$, 50 epochs); TabNet \cite{arik2021tabnet} for 200 epochs (learning rate $2\times10^{-2}$); XGBoost \cite{chen2016xgboost} with default settings. KAN \cite{liu2025kankolmogorovarnoldnetworks} uses $[3,32,c]$ / $[13,64,5]$ with 4 B-spline bases, SGD (learning rate $0.001$). FastKAN \cite{li2024kolmogorovarnoldnetworksradialbasis} shares the KAN configuration with 50 epochs.

\noindent \textbf{UA-KAN} UA-KAN shares the FastKAN configuration with the following additions: learnable RBF width $h$ (initial value $0.001$), RS residual with $\gamma=4$, and the composite loss (Eq.~\ref{eq:proposedloss}) with $\lambda=0.8$, $\tau=0.04$ (SEPSIS) and $\tau=0.1$, learning rate $0.01$ (Heart Disease), trained for 50 epochs on both datasets with architecture $[3,32,c]$ / $[13,32,5]$.


\noindent \textbf{Uncertainty Scoring} OoD detection for KAN-based models uses the per-layer RBF basis entropy (Eq.~\ref{eq:layerentropy}), from the final layer only for UA-KAN and as a training-variance-weighted sum for UA-KAN (weighted) (Eq.~\ref{eq:weighteduq}). FT-Transformer, TabNet, and XGBoost use the predictive entropy of the softmax output:
\begin{equation} \label{eq:entropy}
H(\mathbf{p}) = -\sum_{i=1}^{C} p_i \log p_i,
\end{equation}
where $\mathbf{p} \in \mathbb{R}^C$ is the softmax output over $C$ classes.

\subsection{Experimental Results}



\noindent \textbf{SEPSIS Multi-class (Table~\ref{tab:SEPSIS3class})} UA-KAN achieves the highest AUROC on Wine (0.847), Iris (0.914), and Breast Cancer (0.811). TabPFN leads on the Noised variant (0.712) and achieves the highest in-domain accuracy (55.6\%), yet its OoD performance on structured distributional shifts (Wine, Iris, B.cancer) is among the weakest, suggesting that in-context learning does not naturally generalize to OoD separation on these benchmarks. DUQ shows competitive OoD detection but with high variance. The UA-KAN \textit{weighted} variant underperforms the simple variant across all OoD sets on SEPSIS, attributed to the shallow low-dimensional architecture ($[3, 32, c]$): layer-variance weighting amplifies early-layer contributions that carry limited OoD-discriminative information, diluting the stronger signal at the final layer.

\noindent \textbf{SEPSIS Binary (Table~\ref{tab:SEPSIS2class})} All methods achieve comparable in-domain accuracy ($\approx$75--78\%), with TabPFN reaching the highest (78.2\%). UA-KAN achieves the best AUROC on Wine (0.831), Iris (0.900), and Breast Cancer (0.808), while the weighted variant leads on Noised SEPSIS (0.742). Despite its high accuracy, TabPFN shows poor OoD detection (Wine 0.513, Iris 0.587, B.cancer 0.556), confirming that discriminative performance does not imply reliable uncertainty estimation.

\noindent \textbf{Ablation Study (Table~\ref{tab:ablation})} Each architectural component contributes progressively to OoD performance. Replacing B-splines with RBFs alone (KAN+RBF) does not improve over standard KAN without the RS residual; adding RS yields a substantial gain across all OoD sets. The full uncertainty mechanism produces a further large jump (e.g., Wine: $0.765 \to 0.847$). Adding LogitNorm improves in-domain accuracy ($45.8\% \to 50.2\%$) with a marginal AUROC trade-off, confirming the complementary role of each component.

\noindent \textbf{Computational Efficiency (Table~\ref{tab:flops})} UA-KAN (1{,}765 FLOPs, 0.512\,ms) is competitive with FastKAN and orders of magnitude cheaper than FT-Transformer (3.57M FLOPs). TabPFN incurs an estimated $\approx$6.7\,B FLOPs per inference call due to its in-context learning design, which processes the full training set alongside each test point, resulting in 436\,ms per sample — rendering it impractical for real-time clinical deployment despite strong accuracy.

\noindent \textbf{Heart Disease (Table~\ref{tab:heartdisease})} UA-KAN (weighted) achieves the best AUROC across all OoD sets: Wine (0.728), Breast Cancer (0.746), and Noised (0.908). In contrast to SEPSIS, the weighted variant significantly outperforms the simple variant, as the higher-dimensional input space ($[13,32,5]$) allows intermediate layers to develop more informative representations. TabPFN achieves moderate OoD performance (Noised 0.761) but remains well below UA-KAN weighted, at three orders of magnitude higher inference cost.


\input{tables/ablation.tex}



\input{tables/flops.tex}

\input{tables/heart.tex}

\section{Conclusions}

In this study, we investigated the potential of KANs to serve as trustworthy, uncertainty-aware models. While the compositional and interpretable nature of KANs offers a promising architectural alternative to MLPs, standard KANs lack explicit mechanisms for robust uncertainty quantification and remain prone to overconfident predictions on OoD data. To address these limitations, we proposed a novel UA-KAN architecture that integrates trainable radial basis functions and a custom composite loss. This method strengthens the local representation of epistemic boundaries and directly couples confidence calibration with the functional learning process. Our comprehensive evaluation on clinical datasets demonstrates that the proposed UA-KAN significantly outperforms Deep Ensembles, DUQ and KAN baselines in OoD detection, while maintaining competitive in-distribution accuracy and favorable training efficiency. These results establish that with targeted architectural and optimization-level innovations, KANs can be effectively adapted into reliable, uncertainty-aware systems, positioning them as a practical foundation for safety-critical applications.




\section{Acknowledgment}

This paper has received partial funding from the Hellenic Foundation for Research \& Innovation (H.F.R.I.) scholarship under grant agreement No 20490 (Deep Learning Methodologies for Trustworthy Intelligent Systems). This work was partially supported by the European H2020 project AMBROSIA (no. 101093166). This publication reflects the authors’ views only. The H.F.R.I. and European Commission are not responsible for any use that may be made of the information it contains.

\clearpage
\newpage

\bibliographystyle{IEEEtran}
\bibliography{IEEEabrv,references}


\end{document}

